% Amin Samadi — Resume
% Compile with: tectonic Amin_Samadi_Resume.tex
\documentclass[10pt,letterpaper]{article}

\usepackage[margin=0.65in]{geometry}
\usepackage{xcolor}
\usepackage{titlesec}
\usepackage{enumitem}
\usepackage[hidelinks]{hyperref}

\definecolor{heading}{RGB}{70,100,160}
\urlstyle{same}
\pagestyle{empty}
\setlength{\parindent}{0pt}

\titleformat{\section}{\large\scshape\color{heading}}{}{0em}{}[\vspace{1pt}\titlerule]
\titlespacing*{\section}{0pt}{12pt}{6pt}

\setlist[itemize]{leftmargin=1.2em, itemsep=1pt, parsep=0pt, topsep=3pt, label=\textbullet}

% \entry{title}{dates}{subtitle}{location}
\newcommand{\entry}[4]{%
  \par\textbf{#1}\hfill\textbf{#2}\\
  \textit{#3}\hfill\textit{#4}\par\vspace{1pt}}

\newcommand{\sep}{\,$\cdot$\,}

\begin{document}

\begin{center}
  {\LARGE\bfseries Mohammad Amin Samadi}\\[3pt]
  {\color{heading}Researcher: Learning Sciences, Human--AI Interaction \& Learning Analytics}\\[3pt]
  San Francisco, CA \sep (714) 516-0001 \sep
  \href{mailto:masamadi@uci.edu}{masamadi@uci.edu} \sep
  \href{https://aminsamadi.org}{aminsamadi.org} \sep
  \href{https://github.com/aminsmd}{github.com/aminsmd} \sep
  \href{https://linkedin.com/in/aminsamadi}{linkedin.com/in/aminsamadi}
\end{center}

\section{Summary}

Mixed-methods researcher studying how people learn and build capability with AI. PhD candidate in Education (HCI / Learning Analytics) at UC Irvine; dissertation on human--AI collaboration in teams. I build the instruments and platforms my research questions require: I designed and lead TRAIL, a controlled-study platform for human--AI teaming with randomized assignment, survey instrumentation, and automated data capture. I also ship production LLM systems (currently a Claude-based agent with full evaluation infrastructure at ROI-AI One).

\section{Education}

\entry{Ph.D. Candidate, Education (HCI / Learning Analytics)}{2021 -- Present}
  {University of California, Irvine. Advisor: Nia Nixon}{Irvine, CA}
\begin{itemize}
  \item Dissertation (in progress): human--AI collaboration in teams, built on the TRAIL platform.
  \item UCI School of Education Vandell Dissertation Research Award; Best Short Paper Award, EDM 2025.
\end{itemize}

\entry{B.Sc., Computer Engineering}{2015 -- 2019}
  {Isfahan University of Technology. Advisor: Mohammad Hossein Manshaei}{Isfahan, Iran}

\section{Industry Experience}

\entry{AI Scientist Intern}{Jan 2026 -- Present}
  {ROI-AI One (\href{https://roiaione.com/}{roiaione.com}) -- agentic systems for retail chargeback management}{San Francisco, CA}
\begin{itemize}
  \item Took the Root Cause Analysis tool from concept to production: a multi-agent workflow that ingests case documents, tracks findings in a live ledger, and determines why a chargeback occurred against industry standards; demoed by the CEO to customers.
  \item Built an observability and eval system for agentic workflows: visualizes tool calls, thinking blocks, and agent dispatch, and scores agent performance across multiple methods to auto-flag regressions before production.
  \item Built an email orchestration system that routes inbound messages to the correct workflow, then drafts and sends the outbound response, with every automated action gated behind a per-customer, per-action switch that ships off by default.
  \item Built document intake and portal automation: OCR-based ingestion (AWS Textract + LLM tool-use) paired with automation that logs into retailer portals, imports documents, and submits response packets end-to-end across five retailer portals.
  \item Built an auto-improvement backend that turns internal analyst feedback into a suggested fix, opened automatically as a pull request.
  \item Built the pipeline that turns production agent traces into GRPO/SFT training data (context-window management, Hermes-formatted tool-call data) to distill the agent's reasoning into a smaller, deployable model (Qwen); ran initial training tests to validate the pipeline.
  \item Migrated the team's instance infrastructure from a single Lightsail instance to EC2.
\end{itemize}

\section{Research Experience}

\entry{TRAIL: Team Research and AI Integration Lab}{May 2023 -- Present}
  {Lead researcher and engineer, Nixon Lab, UC Irvine}{Irvine, CA}
\begin{itemize}
  \item Designed and lead TRAIL, a research platform that embeds configurable LLM teammates into real-time, multi-party chatrooms for controlled studies of human--AI collaboration; deployed on AWS and used across multiple lab studies with 50+ concurrent participants.
  \item Built the measurement toolkit the lab's studies run on: a self-hosted survey module, an A/B framework with randomized team assignment and automated data capture, RAG-grounded study materials, and collaborative experiment sharing with role-based permissions and version history.
  \item Led qualitative coding of team discourse data: developed coding schemes, trained and managed coder teams, and established inter-rater reliability before analysis.
  \item Designed and administered survey instruments for studies of team dynamics and perceptions of AI teammates.
  \item Analyzed team interaction data with mixed-effects models, ANOVA, and psycholinguistic methods (LIWC, discourse analysis), linking team discourse dynamics to collaboration outcomes; complemented statistical findings with qualitative analysis of chat transcripts.
  \item Publish public-facing research writeups of TRAIL studies at \href{https://aminsamadi.org/blog}{aminsamadi.org/blog}.
  \item Architected the AI teammate as a multi-agent pipeline (hierarchical memory, Big-Five persona compiler, turn-taking rules calibrated against real chat data) with a simulation/replay tool for testing behavior changes against past experiments before deployment.
  \item Managed a cross-disciplinary team of research interns and MSWE capstone developers.
\end{itemize}

\entry{Multimodal Classroom Video Understanding Pipeline}{Jan 2025 -- Dec 2025}
  {\href{https://github.com/aminsmd/multimodal-transcription}{github.com/aminsmd/multimodal-transcription}, UC Irvine}{Irvine, CA}
\begin{itemize}
  \item Built an AI pipeline that transcribes and semantically segments classroom video to support research on teacher--student interaction quality; multimodal (audio + visual) transcription via Gemini 2.5 Pro, segmentation via temporal embedding similarity, deployed at scale on AWS ECS.
  \item Led the design of a qualitative coding framework for teacher actions, then operationalized it in a multi-agent system that codes classroom video automatically.
  \item Built a human-coded evaluation set from the framework and used it to validate the multi-agent system's coding against expert judgments. Work published in two 2026 papers (see Selected Publications).
\end{itemize}

\section{Skills}

\textbf{Research methods:} mixed-methods research, experimental design, A/B testing, survey design, qualitative coding and inter-rater reliability, instrument development and validation, statistical modeling (mixed-effects models, ANOVA), psycholinguistic analysis (LIWC), discourse analysis, learning analytics, LLM evals and benchmark construction, UX research (HCI)

\smallskip
\textbf{ML \& LLMs:} AI agents and multi-agent systems, Claude Agent SDK workflows, tool use / function calling, RAG, memory systems, prompt engineering, prompt caching, post-training (SFT, GRPO), reward modeling, LLM observability, multimodal LLMs (Claude, Gemini, OpenAI, Qwen)

\smallskip
\textbf{Engineering:} Python, R, TypeScript, SQL; PyTorch, Hugging Face Transformers, TRL, LangChain, CrewAI; Django, Flask, FastAPI, Next.js/React, Socket.IO; AWS (ECS, Fargate, S3, Textract), Docker, Terraform, PostgreSQL, GitHub Actions

\section{Selected Publications}

\begin{itemize}
  \item J. Choi, \textbf{M. A. Samadi}, S. Jaquay, S. Park, N. Nixon. ``Read the Room or Lead the Room: Understanding Socio-Cognitive Dynamics in Human-AI Teaming.'' \textit{Learning Analytics and Knowledge (LAK '26)}.
  \item \textbf{M. A. Samadi}, N. Nixon. ``Personalities at Play: Probing Alignment in AI Teammates.'' Preprint, 2026.
  \item J. Wu, L. Li, \textbf{M. A. Samadi}, et al. ``Validating High-Inference Pedagogical Constructs: Preliminary Evidence from a Multi-agent Architecture for Equity-Focused Mathematics.'' \textit{International Conference on Artificial Intelligence in Education (AIED '26)}.
  \item L. Li, \textbf{M. A. Samadi}, J. Wu, et al. ``Automated Multimodal Transcription for Belonging-Centered Classroom Interaction Analysis: Opportunities and Challenges.'' \textit{International Conference on Artificial Intelligence in Education (AIED '26)}.
  \item \textbf{M. A. Samadi}, S. JaQuay, J. Gu, N. Nixon. ``The AI Collaborator: Bridging Human-AI Interaction in Educational and Professional Settings.'' Preprint, \href{https://arxiv.org/abs/2405.10460}{arxiv.org/abs/2405.10460}.
  \item \textbf{M. A. Samadi}, S. JaQuay, Y. Lin, E. Tajik, S. Park, N. Nixon. ``Minds and Machines Unite: Deciphering Team Social and Cognitive Dynamics with AI.'' \textit{Learning Analytics and Knowledge (LAK '24)}.
  \item \textbf{M. A. Samadi}, D. Shariff, R. Ibarra, N. Nixon. ``Balancing the Scales: Using GPT-4 for Robust Data Augmentation.'' \textit{Companion Proceedings of LAK '25}.
  \item S. Park, D. Shariff, \textbf{M. A. Samadi}, N. Nixon, S. D'Mello. ``From Discourse to Dynamics: Understanding Team Interactions Through Temporally Sensitive NLP.'' \textit{Educational Data Mining (EDM '25)}. \textbf{Best Short Paper Award.}
  \item K. Karbasi, K. Hong, \textbf{M. A. Samadi}, G. Pottie. ``Multi-Agent Collaborative Framework for Math Problem Generation.'' \textit{Educational Data Mining (EDM '25)}.
\end{itemize}

\end{document}
